% Part: sets-functions-relations
% Chapter: sets
% Section: unions-and-intersections

\documentclass[../../../include/open-logic-section]{subfiles}

\begin{document}

\olfileid{sfr}{set}{uni}
\olsection{సమ్మేళనాలు, ఛేదనాలు}

\begin{explain}
\olref[sfr][set][bas]{sec}లో ధర్మనిర్దేశం ద్వారా సమితులను
నిర్వచించాం; అంటే $\Setabs{x}{\phi(x)}$ రూపంలోని నిర్వచనాలను
పరిచయం చేశాం. ఇక్కడ వాడే $\phi$ అనే ధర్మంలో ఇంతకుముందే
నిర్వచించిన సమితులను పేర్కొనవచ్చు. ఉదాహరణకు $A$, $B$ సమితులైతే,
$\Setabs{x}{x \in A \lor x \in B}$ అనే సమితిలో $A$లో గానీ
$B$లో గానీ !!{element}sగా ఉన్న వస్తువులన్నీ ఉంటాయి.
అంటే ఇది $A$, $B$ల \tecase{acc}{!!{element}s} కలిపే సమితి.
దీనిని \olref{fig:union}లో చూడవచ్చు. అందులో గుర్తించిన ప్రాంతం
రెండు సమితులు $A$, $B$లకు చెందిన \tecase{acc}{!!{element}s}
కలిపి చూపిస్తుంది.

\begin{figure}
  \olasset{assets/diagrams/union.tikz}
  \caption{రెండు సమితుల సమ్మేళనం $A \cup B$ అనేది $A$లోని
  \tecase{acc}{!!{element}s}, $B$లోని మూలకాలతో కలిపిన సమితి.}
  \ollabel{fig:union}
\end{figure}

సమితులను కలిపే ఈ ప్రక్రియ ఉపయోగకరమైనదీ, తరచుగా వాడేదీ కనుక
దీనికి ఒక నియతమైన పేరు, ఒక సంకేతం ఇస్తాం.
\end{explain}

\begin{defn}[సమ్మేళనం]
$A$, $B$ అనే రెండు సమితుల \emph{సమ్మేళనం} $A \cup B$:
ఇది $A$లోనో, $B$లోనో, లేదా రెండింటిలోనో !!{element}sగా ఉన్న
వస్తువులన్నింటి సమితి.
\[
A \cup B = \Setabs{x}{x \in A \lor x \in B}
\]
\end{defn}

\begin{ex}
\tecase{acc}{!!{element}s} ఎన్నిసార్లు పేర్కొన్నామన్నది ముఖ్యం కాదు.
అందువల్ల రెండు సమితులకు ఉమ్మడిగా !!a{element} ఉన్నప్పుడు,
వాటి సమ్మేళనంలో ఆ !!{element} ఒక్కసారి మాత్రమే ఉంటుంది.
ఉదాహరణకు $\{ a, b, c\} \cup \{ a, 0, 1\} = \{a, b, c, 0, 1\}$.

ఒక సమితికీ, దాని ఉపసమితికీ సమ్మేళనం పెద్ద సమితే:
$\{a, b, c \} \cup \{a \} = \{a, b, c\}$.

ఏ సమితికైనా శూన్య సమితితో సమ్మేళనం అదే సమితి:
$\{a, b, c \} \cup \emptyset = \{a, b, c \}$.
\end{ex}

\begin{prob}
$A \subseteq B$ అయితే $A \cup B = B$ అని నిరూపించండి.
\end{prob}

\begin{explain}
సమ్మేళనానికి ``ద్వంద్వమైన'' ప్రక్రియను కూడా పరిగణించవచ్చు.
$A$లో !!{element}sగా ఉంటూ, $B$లో కూడా !!{element}sగా ఉండే
\tecase{gen}{!!{element}s} సమితిని ఏర్పరచే ప్రక్రియ ఇది.
దీనిని \emph{ఛేదనం} అంటాం; \olref{fig:intersection}లో దీనిని చూడవచ్చు.
\begin{figure}
  \olasset{assets/diagrams/intersection.tikz}
  \caption{రెండు సమితుల ఛేదనం $A \cap B$ అనేది వాటికి ఉమ్మడిగా
    ఉన్న \tecase{gen}{!!{element}s} సమితి.}
  \ollabel{fig:intersection}
\end{figure}
\end{explain}

\begin{defn}[ఛేదనం]
$A$, $B$ అనే రెండు సమితుల \emph{ఛేదనం} $A \cap B$:
ఇది $A$, $B$ రెండింటిలోనూ !!{element}sగా ఉన్న వస్తువులన్నింటి సమితి.
\[
A \cap B = \Setabs{x}{x \in A \land x \in B}
\]
రెండు సమితుల ఛేదనం శూన్యమైతే వాటిని \emph{వియుక్త} సమితులు అంటాం.
అంటే వాటికి ఉమ్మడి !!{element}s ఏవీ ఉండవు.
\end{defn}

\begin{ex}
రెండు సమితులకు ఉమ్మడి !!{element}s ఏవీ లేకపోతే, వాటి ఛేదనం శూన్యం:
$\{ a, b, c\} \cap \{ 0, 1\} = \emptyset$.

రెండు సమితులకు ఉమ్మడి !!{element}s ఉంటే, వాటన్నింటి సమితే ఛేదనం:
$\{a, b, c \} \cap \{a, b, d \} = \{a, b\}$.

ఒక సమితికీ, దాని ఉపసమితికీ ఛేదనం చిన్న సమితే:
$\{a, b, c\} \cap \{a, b\} = \{a, b\}$.

ఏ సమితికైనా శూన్య సమితితో ఛేదనం శూన్యమే:
$\{a, b, c \} \cap \emptyset = \emptyset$.
\end{ex}

\begin{prob}
$A \subseteq B$ అయితే $A \cap B = A$ అని కచ్చితంగా నిరూపించండి.
\end{prob}

\begin{explain}
రెండుకంటే ఎక్కువ సమితుల సమ్మేళనం లేదా ఛేదనాన్ని కూడా ఏర్పరచవచ్చు.
సాధారణంగా దీనికి ఒక చక్కని పద్ధతి ఉంది. సమ్మేళనం (లేదా ఛేదనం)
ఏర్పరచదలచిన సమితులన్నింటినీ ఒకే సమితిలో చేర్చామనుకుందాం.
అప్పుడు మూల సమితుల సమ్మేళనాన్ని, ఈ సమితిలోని కనీసం ఒక
\tecase{dat}{!!{element}} చెందిన వస్తువులన్నింటి సమితిగా నిర్వచించవచ్చు.
ఛేదనాన్ని, ఈ సమితిలోని ప్రతి \tecase{dat}{!!{element}} చెందిన వస్తువులన్నింటి
సమితిగా నిర్వచించవచ్చు.
\end{explain}

\begin{defn}
$A$ అనేది సమితుల సమితి అయితే, $\bigcup A$ అనేది $A$లోని
\tecase{gen}{!!{element}s} \tecase{gen}{!!{element}s} సమితి:
\begin{align*}
\bigcup A & = \Setabs{x}{x \text{ అనేది } A\text{లోని \tecase{dat}{!!a{element}} చెందుతుంది}},
\text{ అంటే,}\\
& = \Setabs{x}{\text{ఏదో ఒక } B \in A
  \text{ కోసం } x \in B}
\end{align*}
\end{defn}

\begin{defn}
$A$ అనేది సమితుల సమితి అయితే, $\bigcap A$ అనేది $A$లోని
మూలకాలన్నింటికీ ఉమ్మడిగా ఉన్న వస్తువుల సమితి:
\begin{align*}
\bigcap A & = \Setabs{x}{x \text{ అనేది } A\text{లోని ప్రతి \tecase{dat}{!!{element}} చెందుతుంది}},
\text{ అంటే,}\\
 & = \Setabs{x}{\text{ప్రతి } B \in A\text{ కోసం}, x \in B}
\end{align*}
\end{defn}

\begin{ex}
$A = \{ \{ a, b \}, \{ a, d, e \}, \{ a, d \} \}$ అనుకుందాం.
అప్పుడు $\bigcup A = \{ a, b, d, e \}$, $\bigcap A = \{ a \}$.
\end{ex}
\begin{prob}
	$A$ ఒక సమితి అయి, $A \in B$ అయితే $A \subseteq \bigcup B$ అని చూపండి.
\end{prob}

$A_1$, $A_2$, \dots అనే సమితుల క్రమానికీ ఇదే విధంగా చేయవచ్చు:
\begin{align*}
\bigcup_i A_i & = \Setabs{x}{x \text{ అనేది ఏదో ఒక } A_i\text{కు చెందుతుంది}}\\
\bigcap_i A_i & = \Setabs{x}{x \text{ అనేది ప్రతి } A_i\text{కు చెందుతుంది}}.
\end{align*}

సమితులకు ఒక \emph{సూచిక} ఉన్నప్పుడు, అంటే $I$ అనే సమితిలోని
ప్రతి $i \in I$కూ $A_i$ను పరిగణిస్తున్నప్పుడు, ఈ సంక్షిప్త రూపాలను
కూడా వాడవచ్చు:
\begin{align*}
	\bigcup_{i \in I} A_i & = \bigcup \Setabs{A_i }{i \in I}\\
	\bigcap_{i \in I} A_i & = \bigcap\Setabs{A_i}{i \in I}
\end{align*}

చివరగా, $A$లో ఉంటూ $B$లో లేని \tecase{gen}{!!{element}s}
సమితిని పరిశీలిద్దాం. దీనిని \olref{difference}లో చూపించవచ్చు.

\begin{figure}
  \olasset{assets/diagrams/difference.tikz}
  \caption{రెండు సమితుల భేదం $A \setminus B$ అనేది $A$లో
    !!{element}sగా ఉండి, $B$లో !!{element}sగా లేని వస్తువుల సమితి.}
  \ollabel{difference}
\end{figure}

\begin{defn}[భేదం]
\emph{సమితుల భేదం} $A \setminus B$ అనేది $A$లో
!!{element}sగా ఉండి, $B$లో !!{element}sగా లేని వస్తువులన్నింటి సమితి.
అంటే,
\[
A\setminus B = \Setabs{x}{x\in A \text{ మరియు } x \notin B}.
\]
\end{defn}

\begin{prob}
	$A \subsetneq B$ అయితే $B \setminus A \neq \emptyset$ అని నిరూపించండి.
\end{prob}

\end{document}
